% SelInfBench -> IEEE BIBE 2026
% Venue: IEEE BIBE 2026 (EI/IEEE, double-column 10pt, full 8 pages)
% Framing (user-approved 2026-06-19): clinical/reproducibility-first.
% Title keywords: "Medical Imaging" + "Trustworthy Reporting".
% Number discipline: all numbers must trace to paper/BUILD_MAP.md verified csv. Use \todo{} placeholders until verifier-confirmed.

\documentclass[conference]{IEEEtran}
\IEEEoverridecommandlockouts

\usepackage{cite}
\usepackage{amsmath,amssymb,amsfonts}
\usepackage{algorithm}
\usepackage{algorithmic}
\usepackage{graphicx}
\usepackage{textcomp}
\usepackage{xcolor}
\usepackage{booktabs}

% TODO marker (renders red, non-fatal)
\newcommand{\todo}[1]{\textcolor{red}{[TODO: #1]}}

\begin{document}

\title{Are Reported Accuracies in Medical Imaging AI Benchmarks Trustworthy? Diagnosing and Correcting the Winner's Curse via Data Fission}

\author{\IEEEauthorblockN{Anonymous}\IEEEauthorblockA{\todo{author block / affiliation -- fill at camera-ready, keep anonymous if double-blind required}}}

\maketitle

\begin{abstract}
\todo{Abstract -- writer to draft. Cover: (1) medical imaging AI papers report best-of-sweep numbers -> winner's curse inflates reported accuracy and invalidates CIs used for deployment decisions; (2) we measure winner's curse on real medical benchmarks (HAM10000, BraTS); (3) classical CIs undercover the selected model's true accuracy, a data-fission correction restores nominal coverage; (4) we deliver a plug-in corrector. Numbers from BUILD\_MAP only.}
\end{abstract}

\begin{IEEEkeywords}
medical imaging, reproducibility, benchmark, winner's curse, selective inference, data fission, trustworthy reporting
\end{IEEEkeywords}

\section{Introduction}
Clinicians and engineers choosing a model for deployment lean heavily on the
accuracy figures, and the confidence intervals around them, reported in
medical-imaging benchmark studies. Yet the prevailing way those figures are
produced quietly undermines their trustworthiness. A typical study sweeps over
many candidate configurations---combinations of hyperparameters and random
seeds---trains each, and reports the single configuration that attains the best
validation score. Reporting the best of many noisy measurements is precisely the
setting of the \emph{winner's curse}: the reported value is, in expectation,
larger than the true performance of the configuration it singles out, so the
accompanying confidence interval, computed as though that configuration had been
fixed in advance, fails to cover the quantity practitioners care about. The
interval that a clinician reads as ``the model attains $X$ with $95\%$
confidence'' need not contain the deployed model's true accuracy at anything
near its stated rate. This is a reproducibility and trustworthy-reporting
problem at the heart of biomedical engineering: the numbers that drive
model-selection and deployment decisions are reported with an optimistic tilt
and an invalid uncertainty statement.

We stress at the outset that this is a problem of \emph{reporting practice}, not
a claim that the underlying models are flawed. The inflation arises mechanically
from selecting and reporting the maximum over a sweep, regardless of whether the
swept methods are individually sound. Among the many forms of selection that can
distort a benchmark, we target the most common one in day-to-day practice yet
the one least often corrected: reporting the best result of a single method
across its own hyperparameter and seed sweep. Unlike cross-model selection,
where one picks the best among distinct competing methods, this within-method
best-of-sweep reporting is so routine that it is rarely recognized as a
selection event requiring statistical correction at all.

Statistics offers principled remedies for inference after selection---exact
post-selection inference~\cite{lee2016exact}, defenses against the winner's curse
in adaptive analysis~\cite{zrnic2024flexible}, and data
fission~\cite{leiner2023datafission}, which splits each observation into
independent selection and inference channels. These tools, however, have not been
brought to bear on how medical-imaging benchmarks are reported. We close that
gap. Concretely, our contributions are threefold.
\begin{itemize}
  \item \textbf{Framing and method transfer.} To our knowledge this is the first
  work to bring selective inference, and specifically data fission, to the
  reporting habits of medical-imaging benchmarks, recasting best-of-sweep
  reporting as a selection event that admits a conditionally valid correction.
  \item \textbf{Diagnosis and correction on real benchmarks.} On real
  medical-imaging benchmarks (skin-lesion dermoscopy and brain-MRI tumor
  segmentation) we exhibit a measurable winner's curse, show that classical
  confidence intervals undercover the selected configuration's true performance,
  and demonstrate that a data-fission correction restores coverage toward its
  nominal level. We report these effects within the scope of the sweep selection
  we model, not as universal constants over the literature.
  \item \textbf{A plug-in corrector.} We deliver a lightweight report-value
  corrector that attaches to any existing sweep log and returns, for the selected
  configuration, a debiased point estimate together with a conditionally valid
  interval---usable without re-running any training.
\end{itemize}
The remainder of the paper formalizes the setting and the correction
(Section~\ref{sec:method}), reports the empirical diagnosis and recovery
(Section~\ref{sec:results}), and discusses the scope and limitations of the
approach (Section~\ref{sec:discussion}).
\todo{writer added \ref{sec:method} and \ref{sec:discussion}; main owner must add \label{sec:method} under \section{Method} and \label{sec:discussion} under \section{Discussion} (writer not permitted to edit those sections). sec:results already referenced by Method per STORY.}

\section{Related Work}

\paragraph{Selection effects in medical-imaging training.}
The phenomenon we correct has been observed empirically in medical imaging.
\AA{}kesson et al.~\cite{akesson2024random} train a segmentation model across many
random seeds and find that selecting the best-performing seed yields performance
that is statistically significantly higher than that of a large fraction of the
remaining seeds---direct evidence that seed selection manufactures an upward,
statistically significant overestimate. Their study documents the effect but does
not provide a statistical correction for the reported value; the reader is left
with an inflated number and no recipe to recover a valid one. Our work supplies
exactly that missing correction.

\paragraph{Statistical defenses against the winner's curse.}
The winner's curse and, more broadly, inference after data-dependent selection
are well studied in statistics. Exact post-selection
inference~\cite{lee2016exact} conditions on the selection event by characterizing
it as a polyhedral constraint; flexible defenses against the winner's curse have
been developed for adaptive and multi-stage analyses~\cite{zrnic2024flexible};
and data fission~\cite{leiner2023datafission} offers a general device that splits
a single observation into independent selection and inference channels, avoiding
the need for an explicit polyhedral description of the selection rule. These
methods have been applied in settings such as genomics and general machine
learning, but, to our knowledge, none has been instantiated for the reporting of
medical-imaging benchmarks. We are the first to make this transfer, and we pair
it with a corrector that can be applied to existing sweep logs.

\paragraph{Reporting practices in medical-imaging AI.}
Our premise---that best-of-sweep reporting is the norm---rests on documented
reporting habits. Koopmans et al.~\cite{koopmans2025false} review outperformance
claims in medical-imaging AI and find that the large majority report a best
result while only a small minority accompany it with a statistical test. Renard
et al.~\cite{renard2020variability} report that a substantial share of
deep-learning medical-image-segmentation studies omit any measure of variance.
Sculley et al.~\cite{sculley2018winners} raise the same concern of empirical
rigor and the winner's curse for machine learning at large. Together these
sources establish that the dominant reporting habit---report the best number,
seldom report variance, seldom run a statistical test---is consistent with the
winner's curse we study. We are careful not to overstate this evidence: it is
side evidence of under-reported variance and untested claims, together with the
direct seed-selection overestimate of~\cite{akesson2024random}, rather than a
census establishing the precise prevalence of taking the maximum across a sweep.

\section{Method}
\label{sec:method}

\subsection{Problem Setting}
A benchmark report is typically produced by sweeping over $M$ configurations---combinations of hyperparameters and random seeds---training each, and reporting the single configuration that attains the best validation metric. Let $X_i$ denote the observed validation metric for configuration $i \in \{1, \dots, M\}$, modeled as a noisy measurement of an unknown population value $\mu_i$, i.e.\ $X_i = \mu_i + \varepsilon_i$ with $\varepsilon_i \sim \mathcal{N}(0, \sigma^2)$. The selected configuration is
\begin{equation}
i^\star = \arg\max_{i} X_i ,
\end{equation}
and the \emph{naive} report consists of the point value $X_{i^\star}$ together with a standard confidence interval (CI) built around it.

The difficulty is that $i^\star$ is itself a function of the data. Conditioning on the event $\{i^\star = i\}$ shifts the sampling distribution of $X_{i^\star}$ upward: the maximum of $M$ noisy draws systematically overshoots the underlying value of the configuration it selects. Consequently $X_{i^\star}$ is a positively biased estimator of $\mu_{i^\star}$, and a CI centered at $X_{i^\star}$ that ignores the selection event is conditionally invalid---it undercovers $\mu_{i^\star}$, the true performance of the very configuration being reported. This is the winner's curse. Classical post-selection inference~\cite{lee2016exact} addresses it by characterizing the selection event as a polyhedral constraint and conditioning on it; we instead adopt the data-fission construction~\cite{leiner2023datafission}, which decouples selection from inference without requiring an explicit polyhedral characterization of the sweep.

\subsection{Data Fission Construction}
Data fission~\cite{leiner2023datafission} injects independent auxiliary Gaussian noise to split each observation into a \emph{selection channel} and an \emph{inference channel}. For each configuration we draw $Z_i \sim \mathcal{N}(0, \sigma^2)$ independently and form
\begin{equation}
f(X_i) = X_i + \tau Z_i, \qquad
g(X_i) = X_i - \tau^{-1} Z_i ,
\end{equation}
with split parameter $\tau > 0$ ($\tau = 1$ by default). The two channels are constructed so that, under the Gaussian model, $f(X_i)$ and $g(X_i)$ are independent: their covariance $\mathrm{Cov}(f(X_i), g(X_i)) = \sigma^2 - \tau \cdot \tau^{-1} \sigma^2 = 0$, and jointly Gaussian uncorrelated variables are independent. Selection is performed on the selection channel,
\begin{equation}
i^\star = \arg\max_{i} f(X_i) ,
\end{equation}
while inference is performed on the inference channel $g(X_{i^\star})$. Because the selection event $\{i^\star = i\}$ is a function of $f$ alone and is independent of $g$, the conditional distribution of $g(X_{i^\star})$ given the selection is simply its unconditional Gaussian law,
\begin{equation}
g(X_{i^\star}) \;\sim\; \mathcal{N}\!\left(\mu_{i^\star},\; \sigma^2(1 + \tau^{-2})\right) ,
\end{equation}
with no selection-induced shift. A standard (untruncated) interval therefore attains nominal conditional coverage of $\mu_{i^\star}$:
\begin{equation}
\mathrm{CI} = g(X_{i^\star}) \;\pm\; z_{1-\alpha/2}\,\sigma\sqrt{1 + \tau^{-2}} ,
\label{eq:fission-ci}
\end{equation}
where $z_{1-\alpha/2}$ is the upper $\alpha/2$ standard-normal quantile. The point $g(X_{i^\star})$ serves as a debiased estimate of $\mu_{i^\star}$, as the upward selection bias carried by $f$ is removed in the orthogonal inference channel. The validity claim we make is about conditional coverage of the interval in Eq.~\eqref{eq:fission-ci}, established empirically in Section~\ref{sec:results}; we do not interpret interval width or the magnitude of the point adjustment as evidence of validity.

\subsection{Variance Estimation}
The construction requires the noise scale $\sigma$. Because each configuration is trained once, no within-configuration replication is available for a per-configuration bootstrap. We therefore estimate $\sigma$ by the pooled standard deviation of the metric across the $M$ swept configurations,
\begin{equation}
\hat\sigma = \sqrt{\frac{1}{M-1}\sum_{i=1}^{M}\bigl(X_i - \bar X\bigr)^2} ,
\end{equation}
which acts as a conservative proxy for the configuration-to-configuration dispersion that drives the winner's curse. When repeated runs of the same configuration are available, a per-configuration estimate may be supplied directly.

\subsection{Report-Value Corrector}
We package the construction as a tool that attaches to any sweep log. The corrector takes a table of (configuration, metric) rows, estimates $\hat\sigma$ as above (or accepts a user-provided value), and applies the data-fission split with $\tau = 1$. It returns, for the selected configuration, (i) the naive reported value $X_{\arg\max_i X_i}$, (ii) the debiased point estimate $g(X_{i^\star})$, and (iii) the conditionally valid interval of Eq.~\eqref{eq:fission-ci}. The corrector emits neither a width-ratio statistic nor any quantity defined as the difference between the naive value and $g(X_{i^\star})$ as an indicator of validity, since such quantities are not informative about conditional coverage. The procedure is summarized in Algorithm~\ref{alg:corrector}.

\begin{algorithm}[t]
\caption{Data-fission report-value corrector}
\label{alg:corrector}
\begin{algorithmic}[1]
\STATE \textbf{Input:} metrics $X_1,\dots,X_M$; split $\tau$; level $\alpha$; optional $\sigma$
\IF{$\sigma$ not provided}
  \STATE $\hat\sigma \leftarrow$ pooled std of $\{X_i\}$ (ddof $=1$)
\ENDIF
\STATE draw $Z_i \sim \mathcal{N}(0,\hat\sigma^2)$ independently, $i=1,\dots,M$
\STATE $f_i \leftarrow X_i + \tau Z_i$;\quad $g_i \leftarrow X_i - \tau^{-1} Z_i$
\STATE $i^\star \leftarrow \arg\max_i f_i$
\STATE $\hat\sigma_g \leftarrow \hat\sigma\sqrt{1 + \tau^{-2}}$
\STATE \textbf{return} debiased estimate $g_{i^\star}$, interval $g_{i^\star} \pm z_{1-\alpha/2}\,\hat\sigma_g$
\end{algorithmic}
\end{algorithm}

\section{Experiments and Results}
\label{sec:results}
\todo{writer: A1 measurable winner's curse (HAM truthproxy +0.0746, with bootstrap CI from R-experiments) / A2 coverage breakdown + recovery (synthetic coverage\_sim\_v2 + real-data bootstrap coverage) / A3 three benchmarks honest (HAM strong, BraTS ceiling-limited weak, ISIC pre-registered full-test result). All numbers from BUILD\_MAP verified csv.}

\section{Discussion}
\label{sec:discussion}
\todo{writer: scope/limitations -- ISIC low-SNR boundary case (honest), BraTS AUROC ceiling, single-sweep results not universal constants, real-benchmark coverage uses full-test truth proxy (caveat). Corrector as deliverable for the medical imaging community.}

\section{Conclusion}
\todo{writer.}

\bibliographystyle{IEEEtran}
\bibliography{refs}

\end{document}
